\justifying
\NSFCBodyText

\subsubsection{项目特色}

本项目将动态工具任务中的可验证结果与持续学习过程连接起来，用同一套资源口径比较参数适配、轨迹学习、能力保持和行动校准。研究以文本工具沙箱为主体、公开时序任务为补充，既保留 Agent 技术的通用性，也利用申请人在非平稳时序建模方面的积累。四条路线均设置能够独立运行的基础对照，方法增益、额外成本和失效情形分别报告。

\subsubsection{创新点}

\textbf{从轨迹证据的传递关系解释持续适应中的收益与干扰。}已有动态秩、适配器路由和经验重放方法分别解决局部问题。本项目拟研究验证偏差如何经轨迹入选和回放进入低秩更新，并影响历史能力，另行检验更新后的行动置信度变化；通过证据打乱、相关错误注入及容量匹配的组合对照，判断哪些证据信号能够支持有效更新，哪些信号会放大干扰。创新性取决于能否得到超出独立模块组合的可重复机制证据，而非动态秩或闭环结构本身。

\textbf{将执行正确性与训练效用分开建模，研究有限记忆中的轨迹取舍。}轨迹成功标签只能说明一次交互结果，不能直接说明其训练价值。项目拟用独立开发集上的配对探测更新检验候选轨迹的边际作用，在固定生成、验证和训练成本下比较易例、困难正确样本及高置信伪成功样本的选择后果，并考察其通过记忆复用影响旧能力的路径。该设计预期形成可审计的样本取舍依据及其失效条件。

\textbf{围绕策略与工具版本变化建立可核验的行动校准研究。}现有回答级校准不能直接替代动作正确性判断。本项目拟在策略冻结后用独立任务重建行动概率映射，并利用可重置沙箱中的影子结果分析选择偏差，在相同覆盖率与推理预算下区分校准、重规划和拒绝带来的效果。预期回答校准信息何时过期、如何识别过期及恢复需要多少反馈，不将经验风险降低写成任意漂移下的安全保证。
